% Introduction（草稿）
Sports outcome forecasting has attracted sustained attention from
statistics, machine learning, and, more recently, large language
models (LLMs). The task is attractive because outcomes are public,
structured, and frequent, and because a predictive model can be
evaluated not only on accuracy but also on downstream financial
consequences in betting markets. However, the same properties make
football forecasting a demanding test bed: betting markets aggregate
vast amounts of information into closing odds, and any method that
hopes to add value must do so \emph{against} a strong, continuously
updated baseline.

The literature has moved in two directions. On the modeling side,
classical approaches---Poisson and Dixon--Coles models, Elo ratings,
and gradient-boosted trees---provide strong and interpretable
baselines. On the representation side, LLM-based forecasting
systems~\citep{halawi2024approaching, alur2025aia, yang2025llm}
demonstrate that language models can produce calibrated probability
forecasts when given structured evidence. Yet most published
evaluations in the LLM forecasting line measure accuracy and Brier
scores in isolation, without a market baseline, without explicit
leakage controls, and without a betting protocol that would make the
financial claims reproducible.

In this paper, these gaps are addressed with a deliberately conservative
evaluation design. The contributions are three empirical findings,
a protocol, and a public benchmark.

\begin{itemize}
\item \textbf{Finding 1: learned models add no information beyond the
market.} Under a leakage-controlled protocol, no learned method
exceeds de-vigged closing odds (54.2\% accuracy, 0.967 log loss);
model--market disagreement occurs in only 6.3\% of matches and is
equally accurate on both sides (32.9\% each); market features account for 87\% of the aggregate permutation
importance, and referee statistics none. The LLM
matches market-level calibration at negligible cost.
\item \textbf{Finding 2: match difficulty is identifiable ex ante.}
The market's own probability level stratifies accuracy from 40.1\%
to 75.3\%, whereas odds movement and bookmaker dispersion do not
(they are confounded with market strength). A corrected
scenario-complexity score and an uncertainty index stratify matches
monotonically and add roughly ten percentage points beyond market
strength; a no-bet policy on high-uncertainty matches improves
risk-adjusted returns.
\item \textbf{Finding 3: apparent model edge is overconfidence.}
Value-based selection of high-confidence bets loses money (ROI
$-11.5\%$, $-7.1\%$ after calibration); error rates among
high-confidence predictions are higher than the market's (22.3\% vs
16.0\% at confidence $\geq 0.7$); misclassified draws carry higher
confidence than correctly predicted ones; and multi-sample
consistency carries no information. Uncertainty filtering helps;
confidence-based selection hurts.
\item \textbf{Protocol.} All features are pre-match, every statistic
is fit on the training split only, and evaluation is on a held-out
future season (1,104 matches) under an identical betting protocol
(equal-stake and fractional-Kelly, single bookmaker's closing odds),
with bootstrap confidence intervals, coverage reporting, and no-bet
decisions separated from executed bets.
\item \textbf{Benchmark.} The dataset (11,811 top-five-league
matches, 2019/20--2025/26) and the full pipeline are public and
script-driven: every number in this paper is regenerated by
one-command scripts rather than hand-written.
\end{itemize}

The remainder of the paper is organized as follows. Section~\ref{sec:related} reviews
related work. Section~\ref{sec:method} describes the data, features,
models, and risk-control framework. Section~\ref{sec:experiments}
presents the experimental results. Section~\ref{sec:discussion}
discusses implications and limitations, and Section~\ref{sec:conclusion} concludes.
